The two parties to a contract are defined through attributes \attr{CRID} and \attr{CPID}. The first is the party initially \textit{creating} the contract and the second is the counterparty, respectively. Thereby, both \attr{CRID}/\attr{CPID} can take any \textit{role} in the contract or, more specifically, they can be the lender or borrower in a loan (PAM), fixed receiver or payer in an interest rate swap (SWAPS), etc.\\

The \textit{role} of the \attr{CRID} is defined through attribute \attr{CNTRL}. The \textit{role} of \attr{CPID} is derived as the \textit{opposite} side to the contract. Apart from \attr{CNTRL} the attributes are \textit{neutral} to the \textit{role} of \attr{CRID} (or \attr{CPID}).\\

On the other hand, contractual cash flows generated by the POFs and certain states are \textit{role-sensitive}. That is, from the perspective of the \attr{CRID} these quantities represent either claims or obligations. Contract Role Sign function $R$ maps the \attr{CNTRL} attribute into $+1$ indicating a claim or $-1$ indicating an obligation

\[
	R : \attr{CNTRL} \rightarrow \{-1, +1 \}
\]

When multiplying with a cash flow $x$ the Contract Role Sign function thereby defines the direction of that flow:

\begin{itemize}
	\item[$x>0$:] $x$ flows from \attr{CPID} to \attr{CRID}

	\item[$x<0$:] $x$ flows from \attr{CRID} to \attr{CPID}
\end{itemize}

Table \ref{tbl:cntrl} defines the domain of the Contract Role Sign function, i.e. the range of attribute \attr{CNTRL}, with meaning and Contract Role Sign to which the function maps.


% ---------------------- table: contract roles ----------------------
\begin{table}[H]
	\centering
	\begin{tabular}{| p{0.5in}p{1.5in}p{0.2in} |}
	\hline
	\textbf{Value} & \textbf{Meaning} & $\textbf{R}$ \\
	\hline
	RPA & Real position asset & +1 \\
	\hline
	RPL & Real position liability & -1 \\
	\hline
	LG & Long position & +1 \\
	\hline
	ST & Short position & -1 \\
	\hline
